%% 自動生成（phys/gen_phys.py）。直接編集しないこと
\makeatletter
\phdefvol{3}{波動}{ゆれと波を読む}{「なぜ空は青い？」「なぜ音が聞こえる？」に、波で答えよう}
\phdefface{1}{ゆれる}
\phdefface{2}{伝わる}
\phdefface{3}{はかる}
\phdefface{4}{曲がる}
\def\phFrontLine{ビルのゆれも、楽器の音も、光ファイバも。波動の4つの顔で読みとく}
\phdefpart{1}{第1部}{ゆれる：振動と制御}{（波-1〜波-4）}{1}
\phdefpart{2}{第2部}{伝わる：波}{（波-5〜波-8）}{2}
\phdefpart{3}{第3部}{はかる：音・振動の計測と信号}{（波-9〜波-11）}{3}
\phdefpart{4}{第4部}{曲がる：光}{（波-12〜波-14）}{4}
\def\ph@whyA{空は、なぜ青い？　夕焼けは、なぜ赤い？}\def\ph@whyrA{波-13\hspace{0.8em}\bkpageno{228}}
\def\ph@whyB{救急車のサイレンは、なぜ通りすぎると低く聞こえる？}\def\ph@whyrB{波-8\hspace{0.8em}\bkpageno{142}}
\def\ph@whyC{ギターの弦は、なぜおさえる場所で音の高さが変わる？}\def\ph@whyrC{波-6\hspace{0.8em}\bkpageno{106}}
\def\ph@whyD{水に入れたストローは、なぜ折れて見える？}\def\ph@whyrD{波-12\hspace{0.8em}\bkpageno{212}}
\def\ph@whyE{動画の車輪が、逆に回って見えることがあるのはなぜ？}\def\ph@whyrE{波-11\hspace{0.8em}\bkpageno{192}}
\def\ph@whypage{268}
\def\phBackCatch{気になる「なぜ？」から、\\ページをめくってみよう}
\makeatother
